\documentclass[12pt, titlepage]{article}

\usepackage{booktabs}
\usepackage{tabularx}
\usepackage{hyperref}
\usepackage{booktabs}
\usepackage{tabularx}
\usepackage{amssymb}
\usepackage{float}
\hypersetup{
    colorlinks,
    citecolor=blue,
    filecolor=black,
    linkcolor=red,
    urlcolor=blue
}
\usepackage[round]{natbib}

\input{../Comments.text}
\input{../Common.text}

\begin{document}

\title{System Verification and Validation Plan for \progname{}} 
\author{\authname}
\date{\today}
	
\maketitle

\pagenumbering{roman}

\section*{Revision History}

\begin{tabularx}{\textwidth}{p{3cm}p{2cm}X}
\toprule {\bf Date} & {\bf Version} & {\bf Notes}\\
\midrule
Date 1 & 1.0 & Notes\\
Date 2 & 1.1 & Notes\\
\bottomrule
\end{tabularx}

~\\
\wss{The intention of the VnV plan is to increase confidence in the software.
However, this does not mean listing every verification and validation technique
that has ever been devised.  The VnV plan should also be a \textbf{feasible}
plan. Execution of the plan should be possible with the time and team available.
If the full plan cannot be completed during the time available, it can either be
modified to ``fake it'', or a better solution is to add a section describing
what work has been completed and what work is still planned for the future.}

\wss{The VnV plan is typically started after the requirements stage, but before
the design stage.  This means that the sections related to unit testing cannot
initially be completed.  The sections will be filled in after the design stage
is complete.  the final version of the VnV plan should have all sections filled
in.}

\newpage

\tableofcontents

\listoftables
\wss{Remove this section if it isn't needed}

\listoffigures
\wss{Remove this section if it isn't needed}

\newpage

\section{Symbols, Abbreviations, and Acronyms}

\renewcommand{\arraystretch}{1.2}
\begin{tabular}{l l} 
  \toprule		
  \textbf{symbol} & \textbf{description}\\
  \midrule 
  T & Test\\
  \bottomrule
\end{tabular}\\

\wss{symbols, abbreviations, or acronyms --- you can simply reference the SRS
  \citep{SRS} tables, if appropriate}

\wss{Remove this section if it isn't needed}

\newpage

\pagenumbering{arabic}

This document ... \wss{provide an introductory blurb and roadmap of the
  Verification and Validation plan}

\section{General Information}

\subsection{Summary}

% \wss{Say what software is being tested.  Give its name and a brief overview of
%   its general functions.}

The software under test is IDE3.0, a VS Code extension that assists developers
in turning a high-level goal into explicit, reviewable requirements, and keeps
those requirements binding on any AI coding assistant that can implement them.
Without such a tool, the developer struggles with the traceability and
maintainability of goals and constraints since, during development, those
details live only in the chat and prompt phase and are lost when the session
is completed. IDE3.0 instead takes a new spec-driven development approach in
which the important details are treated as artifacts that are versioned in an
intent model, with elicitation artifacts, requirements, rationale, and trace
links that live within the project's repo.

The elicitation follows a staged and structured pipeline, the first stage of
which is the goal decomposition that proposes impacts, deliverables,
and user stories from the developer's goals. A clarification step then asks
about missing rules, exceptions, constraints, and quality needs rather than
the coding assistant inventing them. Finally, a specification stage summarizes
the confirmed information and extracts and classifies explicit requirements.
The developer reviews the proposals at each stage and can return to and edit
earlier stages. From the approved requirements, IDE3.0 generates Java
interfaces with JML contracts that the coding assistant must implement.

IDE3.0 has four components, each covered by its own functional tests:
\begin{itemize}
  \item \textbf{Refinement workflow.} Guides the developer from a
  natural-language goal through decomposition (impact-mapping), clarification, and
  specification to proposed requirements and constraints (FR1--FR3).
  \item \textbf{Review panel and approval gate.} The developer approves,
  edits, or rejects each proposal in VS Code. Code generation is blocked
  until every proposal has a recorded decision (FR4--FR8).
  \item \textbf{Intent model store.} Persists approved artifacts, their
  rationale, and their trace links in the \texttt{.ide3/} folder of the
  developer's workspace (FR9, FR11).
  \item \textbf{Verification arm.} Generates Java interfaces with JML
  contracts from approved requirements and checks the coding assistant's
  implementation against them using OpenJML (FR10, FR13).

\end{itemize}


\subsection{Objectives}

% \wss{State what is intended to be accomplished.  The objective will be around
%   the qualities that are most important for your project.  You might have
%   something like: ``build confidence in the software correctness,''
%   ``demonstrate adequate usability.'' etc.  You won't list all of the qualities,
%   just those that are most important.}

% \wss{You should also list the objectives that are out of scope.  You don't have 
% the resources to do everything, so what will you be leaving out.  For instance, 
% if you are not going to verify the quality of usability, state this.  It is also 
% worthwhile to justify why the objectives are left out.}

% \wss{The objectives are important because they highlight that you are aware of 
% limitations in your resources for verification and validation.  You can't do everything, 
% so what are you going to prioritize?  As an example, if your system depends on an 
% external library, you can explicitly state that you will assume that external library 
% has already been verified by its implementation team.}

The primary objective to be accomplished via this plan is to build confidence
that the IDE3.0 extension correctly carries the developer's intent from a very
high overview of a goal down to verified code generation without loss,
fabrication, or friction. The plan aims to validate the following:

\begin{enumerate}
\item \textbf{Pipeline correctness.} Each elicitation stage produces output that fits and adheres to the shape of the defined intent model schema, and the complete context (original goal, intermediate artifacts, user corrections, confirmed constraints, and unresolved questions) is passed between stages.
\item \textbf{No silent invention.} When a goal or a user story leaves something unspecified, the clarification stage asks the developer for guidance and context rather than inserting an assumption into the requirements.
\item \textbf{Approval gate enforcement.} No artifacts should enter the intent model, nor should any code generation occur, without explicit approval from the human in the loop (and it should be recorded).
\item \textbf{Traceability integrity.} Every approved generated requirement links back to the origination of the goal, user story, and clarification answers, and editing earlier stages correctly ripples through and identifies the affected artifacts downstream.
\item \textbf{Persistence.} The intent model survives across all sessions.
\item \textbf{Contract verification.} Generated JML contracts are well-formed, and OpenJML detects implementations that violate them.
\end{enumerate}

The quality of generated requirements is additionally reviewed against the
characteristics of well-formed requirements in IEEE/ISO/IEC 29148.


\subsection{Custom Documents}

\wss{Summarize the CDs (Custom Documents) that were tackled by this project.  CDs 
can include usability testing, code walkthroughs, user documentation, formal
proof, GenderMag personas, Design Thinking, etc.  CDs should have already
been approved by the course instructor as included in your problem statement.
You can use a pull request to update your CDs (in TeamComposition.csv or
Repos.csv) if your plan changes as a result of the VnV planning exercise.}

\subsection{Relevant Documentation}

\wss{Reference relevant documentation.  This will definitely include your SRS
  and your other project documents (design documents, like MG, MIS, etc).  You
  can include these even before they are written, since by the time the project
  is done, they will be written.  You can create BibTeX entries for your
  documents and within those entries include a hyperlink to the documents.}

\citet{SRS}

\wss{Don't just list the other documents.  You should explain why they are relevant and 
how they relate to your VnV efforts.}

\section{Plan}

% \wss{Introduce this section.  You can provide a roadmap of the sections to
%   come.}

This following section outlines how IDE3.0 will be verified and validated. It defines the verification and validation team and each members responsibilities, then describes the approaches that will be 
taken to verify SRS, the design, this plan, and the implementation. It also lists the automated testing and verification tools, and describes how the software will be validated with our supervisor. 


\subsection{Verification and Validation Team}

% \wss{Your teammates.  Maybe your supervisor.
%   You should do more than list names.  You should say what each person's role is
%   for the project's verification.  A table is a good way to summarize this information.}

\begin{table}[H]
  \centering
  \renewcommand{\arraystretch}{1.3}
  \begin{tabularx}{\textwidth}{|p{3cm}|p{3.2cm}|X|}
  \hline
  \textbf{Name} & \textbf{Role} & \textbf{Responsibilities} \\
  \hline
  Qusay & Technical Verification Co-Lead &
  Develops and executes automated tests for the elicitation pipeline,
  approval gate, intent model store, and verification arm. Maintains the
  CI pipeline and test coverage reporting. Leads design reviews and code
  walkthroughs. \\
  \hline
  Rishab & Technical Verification Co-Lead &
  Shares all responsibilities with Qusay. The two leads split test
  ownership by component and review each other's tests, so no component
  is verified only by its author. \\
  \hline
  [Name] & [Role] & [Responsibilities] \\
  \hline
  [Name] & [Role] & [Responsibilities] \\
  \hline
  [Name] & [Role] & [Responsibilities] \\
  \hline
  Dr.\ Istvan David & Supervisor &
  Provides expert review of the SRS and design in structured review
  meetings, with a focus on the elicitation workflow and traceability
  model. \\
  \hline
  \end{tabularx}
  \caption{Verification and validation team}
  \label{tab:vnv-team}
  \end{table}

\subsection{SRS Verification}

\wss{List any approaches you intend to use for SRS verification.  This may
  include ad hoc feedback from reviewers, like your classmates (like your
  primary reviewer), or you may plan for something more rigorous/systematic.}

\wss{If you have a supervisor for the project, you shouldn't just say they will
read over the SRS.  You should explain your structured approach to the review.
Will you have a meeting?  What will you present?  What questions will you ask?
Will you give them instructions for a task-based inspection?  Will you use your
issue tracker?}

\wss{Maybe create an SRS checklist?}

\subsection{Design Verification}

% \wss{Plans for design verification}

% \wss{The review will include reviews by your classmates}

% \wss{Create a checklists?}

\paragraph{Peer review.}
Team members and classmates will review the design documents using the
checklist in Table~\ref{tab:design-checklist}. Reviewers will document
any criteria that are not satisfied and suggest improvements. All
feedback will be tracked as GitHub Issues and addressed in design
refinement meetings.

\paragraph{Traceability check.}
The technical verification leads will maintain a traceability matrix
mapping each requirement in the SRS to the module(s) that satisfy it.
Any requirement without a module, or any module that satisfies no
requirement, will be flagged and resolved.

\paragraph{Supervisor review.}
A formal review with Dr.\ David will verify three critical aspects of the
design: (1) the elicitation pipeline, including the state passed between
the goal decomposition, clarification, and specification stages (2) the
intent model schema and its trace links and (3) the approval gate and
the generation of Java interfaces with JML contracts. We will prepare a
module diagram, the intent model schema, and a UML sequence diagram
tracing one example goal end-to-end through the system. Action items
from the meeting will be tracked as GitHub Issues.

\begin{table}[H]
\centering
\renewcommand{\arraystretch}{1.3}
\begin{tabularx}{\textwidth}{|>{\raggedright\arraybackslash}p{3.5cm}|>{\raggedright\arraybackslash}X|}
\hline
\textbf{Focus Area} & \textbf{Verification Tasks} \\
\hline
Core Architecture &
$\square$ Every SRS requirement maps to at least one module \newline
$\square$ Modules are small, individually testable, and hide a single
design decision \newline
$\square$ External dependencies (LLM provider, orchestration framework,
OpenJML) are isolated behind interfaces \\
\hline
Elicitation Pipeline &
$\square$ Full context is passed between stages, not only the final user
stories \newline
$\square$ Editing an earlier stage has a defined effect on downstream
artifacts \\
\hline
Intent Model &
$\square$ Schema represents every artifact type and its trace links
\newline
$\square$ Storage format in \texttt{.ide3/} is versionable with the
repository \\
\hline
IDE Integration &
$\square$ VS Code extension API usage reviewed \newline
$\square$ No path allows an artifact to bypass the approval gate \\
\hline
Verification Arm &
$\square$ Mapping from approved requirements to JML contracts is
defined \\
\hline
General &
$\square$ Terminology is consistent with the SRS \\
\hline
\end{tabularx}
\caption{Design verification checklist}
\label{tab:design-checklist}
\end{table}

\subsection{Verification and Validation Plan Verification}

% \wss{The verification and validation plan is an artifact that should also be
% verified.  Techniques for this include review and mutation testing.}

% \wss{The review will include reviews by your classmates}

% \wss{Create a checklists?}

The VnV plan is a project artifact that must itself be verified for
completeness, feasibility, and effectiveness. It will be verified
through team, supervisor, and peer reviews, a feasibility assessment,
mutation testing, and validation of the rubric used to evaluate
LLM-generated artifacts.

\paragraph{Team Review.}
All team members will review the plan using the checklist, confirming that
each objective in Section 2.2 is covered by at least one test and that each
test case is executed as written. The plan will also be checked for
consistency with the SRS and design documents, including component names,
requirement numbers, and tools.

\paragraph{Supervisor Review.}
Dr.\ David will review the test plan for the elicitation pipeline and the
verification arm. We will ask whether the planned tests would give the
developers confidence, or whether there are any crucial testing gaps that
need to be addressed. We will further ask him: Would the failures and the
successes give a more than adequate guarantee that the components function
as they are intended to?

\paragraph{Peer Review.}
There will be a peer review of our plan that is dictated by this document lifecycle, and the issues and recommendations for changes in our approach will be documented on GitHub under Issues.

\paragraph{Feasibility Assessment.}
The technical verification co-leads will confirm that every planned activity is realistic within the project timeline and the team's current skills, time budget, and ability to overcome a not so steep learning curve (if it comes down to it). Activities that cannot be completed in time will be rescoped or listed as out of scope, rather than left unexecuted.

\paragraph{Mutation Testing.}
To verify that the planned tests can detect defects, there will be an intentional introduction of faults injected into the implementation to see if the test cases fail. Any fault that no test detects shows there is a gap in the plan, and a test case will be added to cover it.

\paragraph{Evaluation Rubric Validation.}
The quality of the LLM-generated artifacts will be evaluated against an evaluation (eval) suite. We will manually score a set of deliberately flawed outputs (such as requirements having invented constraints or missing trace links) and confirm that each one falls below the pass threshold. If any of the flawed outputs are considered acceptable, then the evals and acceptance rubric would need to be revised. When the LLM-as-judge is used, its scores will be compared with our manual ratings on a sample of acceptable outputs.

\paragraph{Tracking and Success Metrics.}
All findings will be logged in GitHub Issues and addressed in our team meetings, so that the plan can be refined and retuned. Progress can be measured via requirement coverage, mutation detection rate, and finally issue resolution rate.

\begin{table}[H]
  \centering
  \renewcommand{\arraystretch}{1.3}
  \begin{tabularx}{\textwidth}{|>{\raggedright\arraybackslash}p{3.5cm}|>{\raggedright\arraybackslash}X|}
  \hline
  \textbf{Criteria} & \textbf{Verification Tasks} \\
  \hline
  Coverage &
  $\square$ Every objective in Section~2.2 is covered by at least one
  test \newline
  $\square$ Every functional and non-functional requirement traces to at
  least one test case \newline
  $\square$ Every component (refinement workflow, review panel and
  approval gate, intent model store, verification arm) has dedicated
  tests \\
  \hline
  Test Quality &
  $\square$ Each test has a defined initial state, input, and expected
  output \newline
  $\square$ Tests of LLM-driven stages check structure and behaviour,
  not exact generated wording \\
  \hline
  Effectiveness &
  $\square$ Each seeded fault is detected by at least one test
  \newline
  $\square$ Each deliberately flawed output scores below the rubric
  threshold \newline
  $\square$ LLM-as-judge scores agree with manual ratings on the sampled
  outputs \\
  \hline
  Feasibility &
  $\square$ All planned activities fit within the project timeline \\ 
  \hline
  Document Quality &
  $\square$ The plan is consistent with the SRS and design documents
  \newline
  $\square$ References are complete and correct \newline
  $\square$ The document is free of spelling, grammatical, and
  formatting errors \\
  \hline
  Process &
  $\square$ All review feedback is tracked as GitHub Issues \newline
  $\square$ Traceability tables are updated when requirements change \\
  \hline
  \end{tabularx}
  \caption{VnV plan verification checklist}
  \label{tab:vnv-checklist}
  \end{table}


\subsection{Implementation Verification}

\wss{You should at least point to the tests listed in this document and the unit
  testing plan.}

\wss{In this section you would also give any details of any plans for static
  verification of the implementation.  Potential techniques include code
  walkthroughs, code inspection, static analyzers, etc.}

\wss{The final class presentation in CAS 741 could be used as a code
walkthrough.  There is also a possibility of using the final presentation (in
CAS741) for a partial usability survey.}

\subsection{Automated Testing and Verification Tools}

\wss{What tools are you using for automated testing.  Likely a unit testing
  framework and maybe a profiling tool, like ValGrind.  Other possible tools
  include a static analyzer, make, continuous integration tools, test coverage
  tools, etc.  Explain your plans for summarizing code coverage metrics.
  Linters are another important class of tools.  For the programming language
  you select, you should look at the available linters.  There may also be tools
  that verify that coding standards have been respected, like flake9 for
  Python.}

\wss{If you have already done this in the development plan, you can point to
that document.}

\wss{The details of this section will likely evolve as you get closer to the
  implementation.}

\subsection{Software Validation}

\wss{If there is any external data that can be used for validation, you should
  point to it here.  If there are no plans for validation, you should state that
  here.}

\wss{You might want to use review sessions with the stakeholder to check that
the requirements document captures the right requirements.  Maybe task based
inspection?}

\wss{For those capstone teams with an external supervisor, the Rev 0 demo should 
be used as an opportunity to validate the requirements.  You should plan on 
demonstrating your project to your supervisor shortly after the scheduled Rev 0 demo.  
The feedback from your supervisor will be very useful for improving your project.}

\wss{For teams without an external supervisor, user testing can serve the same purpose 
as a Rev 0 demo for the supervisor.}

\wss{This section might reference back to the SRS verification section.}

\section{System Tests}

\wss{There should be text between all headings, even if it is just a roadmap of
the contents of the subsections.}

\subsection{Tests for Functional Requirements}

\wss{Subsets of the tests may be in related, so this section is divided into
  different areas.  If there are no identifiable subsets for the tests, this
  level of document structure can be removed.}

\wss{Include a blurb here to explain why the subsections below
  cover the requirements.  References to the SRS would be good here.}

\subsubsection{Requirement Review and Approval}

This area tests the presentation, approval, rejection, and revision of
engineering proposals required by FR4-FR6 in the SRS \citep{SRS}.
Related checks cover draft availability (FR3), iterative refinement (FR2),
recorded review decisions (FR8), and persistence of approved intent (FR9).
Test identifiers use \texttt{RRA}, a group code, and a case number:
\texttt{PR} for presentation, \texttt{AD} for accept/decline,
\texttt{ED} for editing through refinement, and \texttt{PA} for pending item
amendment and removal. Accept and decline correspond to approve and reject
in the SRS.

All cases are manual system tests performed through the extension's review
workflow. Each case uses a separate proposal unless stated otherwise.
Persistence checks reopen the same project in a new session and retrieve
its recorded artifacts and decisions, without assuming a storage location
or mechanism. Rejected proposals may be retained as history, but must not
be recorded as approved intent. For generated revisions, the tester checks
the specified meaning rather than exact wording. Report related requirements
below are sample input to \progname{}, not requirements for \progname{} itself.

\paragraph{Presentation}

\begin{enumerate}

\item{RRA-PR-1: Present Requirements and Constraints\\}

Requirements: FR3, FR4.

Control: Manual.

Initial State: A pending proposal contains the requirement ``Users shall be
able to export reports and the constraint ``Exports shall use CSV format.

Input: The engineer opens the proposal for review.

Output: Both items are available for inspection with their complete content
and pending review status. Neither is presented as approved.

Test Case Derivation: FR4 requires requirements and quality constraints to be
presented for review; FR3 requires the draft to be available.

How test will be performed: The tester will open the prepared proposal and
compare the displayed items and review status with the known draft content.

\item{RRA-PR-2: Present an Implementation Contract\\}

Requirements: FR4.

Control: Manual.

Initial State: A pending contract specifies a withdrawal operation with amount
greater than zero, amount no greater than the current balance, and a resulting
balance equal to the previous balance minus the amount.

Input: The engineer opens the contract for review.

Output: The operation, both preconditions, and the postcondition are available
for inspection. The contract is identified as pending review.

Test Case Derivation: FR4 requires contracts to be presented before
corresponding code generation. The engineer must be able to inspect the
proposed conditions.

How test will be performed: The tester will compare the displayed contract
with the prepared draft, checking that the precondition and the
postcondition are not omitted or altered.

\item{RRA-PR-3: Distinguish Multiple Proposal Items\\}

Requirements: FR4, FR6.

Control: Manual.

Initial State: A pending proposal contains requirements for report export and
report preview, plus a CSV export constraint.

Input: The engineer inspects each item and selects the report preview
requirement for amendment.

Output: All three items can be inspected and distinguished. The selected
content corresponds to report preview, with no change to any item.

Test Case Derivation: FR4 requires presentation of the proposal, and FR6
requires individual items to be identifiable for amendment or removal.

How test will be performed: The tester will inspect all three items, select
the preview requirement, and verify that the selected content matches that
item.

\end{enumerate}

\paragraph{Accept and Decline}

\begin{enumerate}

\item{RRA-AD-1: Approve a Requirements Proposal\\}

Requirements: FR5, FR8, FR9.

Control: Manual.

Initial State: A pending proposal contains a report export requirement and a
CSV export constraint.

Input: The engineer approves the proposal.

Output: The approval is recorded for the reviewed content. Both approved items
are retrievable unchanged after reopening the project in a new session.

Test Case Derivation: FR5 permits approval, FR8 requires the decision to be
recorded internally, and FR9 requires approved requirements and constraints to
be retained in the intent model.

How test will be performed: The tester will approve the proposal, check its
decision and content, then reopen the project in a new session and retrieve
the approved items.

\item{RRA-AD-2: Approve a Contract\\}

Requirements: FR4, FR5, FR8.

Control: Manual.

Initial State: A pending withdrawal contract contains the two preconditions
and postcondition described in RRA-PR-2.

Input: The engineer reviews and approves the contract.

Output: The approval is associated with the reviewed contract and its
conditions. The contract and its approval remain retrievable after reopening
the project in a new session.

Test Case Derivation: FR4 includes contracts in the review workflow, while FR5
and FR8 require approval and recording of the decision.

How test will be performed: The tester will inspect and approve the contract,
reopen the project in a new session, and compare the retrieved contract and
decision with the reviewed version.

\item{RRA-AD-3: Reject a Requirements Proposal\\}

Requirements: FR5, FR8, FR9.

Control: Manual.

Initial State: A report export proposal is pending. A separate report preview
requirement has already been approved.

Input: The engineer rejects the pending export proposal.

Output: The rejection remains recorded after reopening the project. The
rejected items do not appear as approved intent, and the previously approved
preview requirement remains unchanged.

Test Case Derivation: FR5 permits rejection and FR8 requires recording the
decision. FR9 distinguishes approved intent from a rejected proposal.

How test will be performed: The tester will reject the export proposal, reopen
the project in a new session, and inspect the recorded decision and approved
artifacts. Retention of the rejected draft as history is acceptable.

\item{RRA-AD-4: Reject a Contract\\}

Requirements: FR4, FR5, FR8.

Control: Manual.

Initial State: A pending withdrawal contract incorrectly permits a negative
withdrawal amount. An unrelated contract is already approved.

Input: The engineer rejects the pending withdrawal contract.

Output: The rejection remains recorded after reopening the project. The
rejected contract is not marked approved, and the unrelated approved contract
is unchanged.

Test Case Derivation: FR4 includes contracts in review, and FR5 and FR8
require rejection to be supported and recorded without assigning approval.

How test will be performed: The tester will reject the contract, reopen the
project in a new session, and check the rejected contract's decision and the
unrelated approved contract.

\end{enumerate}

\paragraph{Editing Through Refinement}

\begin{enumerate}

\item{RRA-ED-1: Add Detail to a Requirement\\}

Requirements: FR2, FR3, FR5, FR6.

Control: Manual.

Initial State: A pending proposal states ``Users shall be able to export reports.

Input: The engineer edits the proposal to specify CSV exports with column headings.

Output: The revised proposal retains report export and includes both CSV
format and column headings. It remains pending and available for review.

Test Case Derivation: FR5 and FR6 permit edits; FR2 requires refinement to
incorporate further information from the engineer.

How test will be performed: The tester will submit the edit, complete any
requested clarification, and inspect the revised proposal for both added
details and its pending status.

\item{RRA-ED-2: Correct a Contract Condition\\}

Requirements: FR2, FR4, FR5.

Control: Manual.

Initial State: A pending withdrawal contract requires amount greater than zero
but does not limit the amount to the available balance.

Input: The engineer edits the proposal to state that overdrafts are prohibited
and the withdrawal amount must not exceed the current balance.

Output: The revised contract includes both a positive-amount condition and a
balance limit. Its result remains consistent with subtracting the withdrawal
amount, and the revised contract is pending review.

Test Case Derivation: FR4 includes contracts in review, FR5 permits editing,
and FR2 requires the refinement to reflect the clarified constraint.

How test will be performed: The tester will submit the correction, complete
any requested clarification, and inspect the resulting conditions and review
status.

\item{RRA-ED-3: Replace a Misunderstood Requirement\\}

Requirements: FR2, FR3, FR5, FR6.

Control: Manual.

Initial State: A pending proposal requires reports to be emailed automatically
to all registered users every day.

Input: The engineer replaces that intent with on-demand report downloads and
explicitly prohibits automatic report emails.

Output: The revised proposal supports on-demand downloads and prohibits
automatic report emails. Daily scheduling and delivery to all users are not
retained as confirmed requirements. The revision remains pending review.

Test Case Derivation: FR5 and FR6 allow correction of a pending proposal,
including substantial changes. FR2 requires subsequent refinement to follow
the revised intent.

How test will be performed: The tester will replace the requirement, complete
any requested clarification, and inspect the resulting proposal for the new
behaviour and absence of the contradicted assumptions.

\item{RRA-ED-4: Retain Successive Edits\\}

Requirements: FR2, FR3, FR5, FR6.

Control: Manual.

Initial State: A pending proposal requires report exports in CSV format.

Input: The engineer first adds column headings, then makes a second edit
requiring export of only the records selected by the user.

Output: The resulting proposal includes CSV format, column headings, and
selected record filtering. It remains pending and retains these details after
reopening the project in a new session.

Test Case Derivation: FR2 requires iterative interaction and FR3 requires
draft retention. Successive edits under FR5 and FR6 must preserve compatible
information from earlier revisions.

How test will be performed: The tester will submit the two edits in sequence,
inspect the resulting proposal, then reopen the project in a new session and
retrieve the pending revision.

\end{enumerate}

\paragraph{Amendment and Removal of Pending Items}

\begin{enumerate}

\item{RRA-PA-1: Amend an Individual Requirement\\}

Requirements: FR3, FR6.

Control: Manual.

Initial State: A pending proposal contains separate requirements for report
export and report preview, plus a CSV export constraint.

Input: The engineer amends only the export requirement to export records
selected by the user.

Output: The export requirement reflects the amendment. The preview requirement
and CSV constraint are unchanged. The saved proposal remains pending after
reopening the project.

Test Case Derivation: FR6 permits amendment of an individual pending
requirement; FR3 requires the revised draft to remain available.

How test will be performed: The tester will amend the export item, reopen the
project in a new session, and compare each retrieved item and its review
status with the expected proposal.

\item{RRA-PA-2: Amend an Individual Quality Constraint\\}

Requirements: FR3, FR6.

Control: Manual.

Initial State: A pending proposal contains report export and report preview
requirements, with a separate constraint that exported reports be readable
using assistive technology.

Input: The engineer amends only the constraint to specify that exported column
headings must be available to screen readers.

Output: The constraint reflects the revised accessibility condition. Both
functional requirements are unchanged, and the revised proposal remains
pending after reopening the project.

Test Case Derivation: FR6 permits amendment of an individual quality
constraint; FR3 requires retention of the revised draft.

How test will be performed: The tester will amend the constraint, reopen the
project in a new session, and inspect the retrieved constraint, other items,
and pending review status.

\item{RRA-PA-3: Remove an Individual Requirement\\}

Requirements: FR3, FR6.

Control: Manual.

Initial State: A pending proposal contains report export and report preview
requirements, plus a CSV constraint associated with export.

Input: The engineer removes only the report preview requirement.

Output: The preview requirement is absent from the pending proposal. The
export requirement and CSV constraint remain unchanged, and the removal
persists after reopening the project.

Test Case Derivation: FR6 permits removal of an individual pending
requirement, while FR3 requires the resulting draft to remain available for
review.

How test will be performed: The tester will remove the preview item, reopen
the project in a new session, and inspect the pending proposal for the removed
item and the unchanged remaining items.

\item{RRA-PA-4: Remove an Individual Quality Constraint\\}

Requirements: FR3, FR6.

Control: Manual.

Initial State: A pending proposal contains report export and report preview
requirements, plus a separate accessibility constraint for exported reports.

Input: The engineer removes only the accessibility constraint.

Output: The constraint is absent from the pending proposal. Both requirements
remain unchanged, and the proposal remains pending with the removal preserved
after reopening the project.

Test Case Derivation: FR6 permits removal of an individual quality constraint
without removing the other pending items. FR3 requires the revised draft to
remain available.

How test will be performed: The tester will remove the constraint, reopen the
project in a new session, and check that the retrieved pending proposal
contains both requirements without the removed constraint.

\end{enumerate}

\subsubsection{Area of Testing2}

...

\subsection{Tests for Nonfunctional Requirements}

\wss{The nonfunctional requirements for accuracy will likely just reference the
  appropriate functional tests from above.  The test cases should mention
  reporting the relative error for these tests.  Not all projects will
  necessarily have nonfunctional requirements related to accuracy.}

\wss{For some nonfunctional tests, you won't be setting a target threshold for
passing the test, but rather describing the experiment you will do to measure
the quality for different inputs.  For instance, you could measure speed versus
the problem size.  The output of the test isn't pass/fail, but rather a summary
table or graph.}

\wss{Tests related to usability could include conducting a usability test and
  survey.  The survey will be in the Appendix.}

\wss{Static tests, review, inspections, and walkthroughs, will not follow the
format for the tests given below.}

\wss{If you introduce static tests in your plan, you need to provide details.
How will they be done?  In cases like code (or document) walkthroughs, who will
be involved? Be specific.}

\subsubsection{Area of Testing1}
		
\paragraph{Title for Test}

\begin{enumerate}

\item{test-id1\\}

Type: Functional, Dynamic, Manual, Static etc.
					
Initial State: 
					
Input/Condition: 
					
Output/Result: 
					
How test will be performed: 
					
\item{test-id2\\}

Type: Functional, Dynamic, Manual, Static etc.
					
Initial State: 
					
Input: 
					
Output: 
					
How test will be performed: 

\end{enumerate}

\subsubsection{Area of Testing2}

...

\subsection{Traceability Between Test Cases and Requirements}

\wss{Provide a table that shows which test cases are supporting which
  requirements.}

\section{Unit Test Description}

\wss{This section should not be filled in until after the MIS (detailed design
  document) has been completed.}

\wss{Reference your MIS (detailed design document) and explain your overall
philosophy for test case selection.}  

\wss{To save space and time, it may be an option to provide less detail in this section.  
For the unit tests you can potentially layout your testing strategy here.  That is, you 
can explain how tests will be selected for each module.  For instance, your test building 
approach could be test cases for each access program, including one test for normal behaviour 
and as many tests as needed for edge cases.  Rather than create the details of the input 
and output here, you could point to the unit testing code.  For this to work, you code 
needs to be well-documented, with meaningful names for all of the tests.}

\subsection{Unit Testing Scope}

\wss{What modules are outside of the scope.  If there are modules that are
  developed by someone else, then you would say here if you aren't planning on
  verifying them.  There may also be modules that are part of your software, but
  have a lower priority for verification than others.  If this is the case,
  explain your rationale for the ranking of module importance.}

\subsection{Tests for Functional Requirements}

\wss{Most of the verification will be through automated unit testing.  If
  appropriate specific modules can be verified by a non-testing based
  technique.  That can also be documented in this section.}

\subsubsection{Module 1}

\wss{Include a blurb here to explain why the subsections below cover the module.
  References to the MIS would be good.  You will want tests from a black box
  perspective and from a white box perspective.  Explain to the reader how the
  tests were selected.}

\begin{enumerate}

\item{test-id1\\}

Type: \wss{Functional, Dynamic, Manual, Automatic, Static etc. Most will
  be automatic}
					
Initial State: 
					
Input: 
					
Output: \wss{The expected result for the given inputs}

Test Case Derivation: \wss{Justify the expected value given in the Output field}

How test will be performed: 
					
\item{test-id2\\}

Type: \wss{Functional, Dynamic, Manual, Automatic, Static etc. Most will
  be automatic}
					
Initial State: 
					
Input: 
					
Output: \wss{The expected result for the given inputs}

Test Case Derivation: \wss{Justify the expected value given in the Output field}

How test will be performed: 

\item{...\\}
    
\end{enumerate}

\subsubsection{Module 2}

...

\subsection{Tests for Nonfunctional Requirements}

\wss{If there is a module that needs to be independently assessed for
  performance, those test cases can go here.  In some projects, planning for
  nonfunctional tests of units will not be that relevant.}

\wss{These tests may involve collecting performance data from previously
  mentioned functional tests.}

\subsubsection{Module ?}
		
\begin{enumerate}

\item{test-id1\\}

Type: \wss{Functional, Dynamic, Manual, Automatic, Static etc. Most will
  be automatic}
					
Initial State: 
					
Input/Condition: 
					
Output/Result: 
					
How test will be performed: 
					
\item{test-id2\\}

Type: Functional, Dynamic, Manual, Static etc.
					
Initial State: 
					
Input: 
					
Output: 
					
How test will be performed: 

\end{enumerate}

\subsubsection{Module ?}

...

\subsection{Traceability Between Test Cases and Modules}

\wss{Provide evidence that all modules have been considered.}
				
\bibliographystyle{plainnat}

\bibliography{../../refs/References}

\newpage

\section{Appendix}

This is where you can place additional information.

\subsection{Symbolic Parameters}

The definition of the test cases will call for SYMBOLIC\_CONSTANTS.
Their values are defined in this section for easy maintenance.

\subsection{Usability Survey Questions?}

\wss{This is a section that would be appropriate for some projects.}

\newpage{}

\section{Appendix --- Generative AI Use Disclosure}

\wss{ChatGPT (version 5) was used to brainstorm ideas for the template for this disclose section.}

\subsection{AI Tools Used}

\wss{Give the name of the AI tool(s) used and the version.}

\subsection{How and Where Used}

\wss{Explain what the tool(s) were used for.  Examples include: brainstorming,
explaining a technical concept, reviewing the document, generating or modifying
text, generating or modifying diagrams, identifying errors or omissions or
inconsistencies.}

\wss{For each of the uses, identify which parts of the document were affected.}

\wss{You do not need to report every interaction, but you should disclose the 
uses that materially contributed to your work}

\subsection{Verification of AI-Generated Content}

\wss{\begin{itemize}
    \item Describe how the team checked AI-generated or AI-assisted material
    before including it in the document.
    \item In particular, verify factual claims, technical assertions,
    requirements, calculations, references, and other information for which an
    AI system may produce incorrect or fabricated results.
    \item \textbf{The team is responsible for the accuracy, completeness,
consistency, and appropriateness of the submitted document, regardless of
whether material was generated by a human or an AI system.}
\end{itemize}}

\newpage{}

\section{Appendix --- Reflection}

\wss{This section is not required for CAS 741}

The information in this section will be used to evaluate the team members on the
graduate attribute of Lifelong Learning.

\input{../Reflection.text}

\begin{enumerate}
  \item What went well while writing this deliverable? 
  \item What pain points did you experience during this deliverable, and how
    did you resolve them?
  \item What knowledge and skills will the team collectively need to acquire to
  successfully complete the verification and validation of your project?
  Examples of possible knowledge and skills include dynamic testing knowledge,
  static testing knowledge, specific tool usage, Valgrind etc.  You should look to
  identify at least one item for each team member.
  \item For each of the knowledge areas and skills identified in the previous
  question, what are at least two approaches to acquiring the knowledge or
  mastering the skill?  Of the identified approaches, which will each team
  member pursue, and why did they make this choice?
\end{enumerate}

\end{document}
